%=======================================================================
%  3  BIO-INSPIRED DESIGN PROCESS
%  Template: "Document the research and design path explicitly. The report
%  should make the reasoning chain from challenge to biological mechanism,
%  technical concept, validation and iteration visible."
%=======================================================================
\section{Bio-Inspired Design Process}

The chain this section has to make visible is
challenge $\rightarrow$ function $\rightarrow$ biology $\rightarrow$ abstraction
$\rightarrow$ concept $\rightarrow$ engineering check $\rightarrow$ evaluation
$\rightarrow$ iteration. Section~2 delivered the first two links. This section
delivers the biology and the abstraction, and hands a set of transferable
principles to the concept development in Section~\ref{sec:creating}.

%-----------------------------------------------------------------------
\subsection{Discovering \& Abstracting}
\label{sec:discovering}

\subsubsection{Research path}
\label{sec:researchpath}

Every search started from a verb in the function list of
Table~\ref{tab:functions}, never from an organism. This follows the rule that
biological content is organised by the function it serves, which is also the
structure of the Biomimicry Taxonomy and therefore of AskNature
\parencite{asknature}. Table~\ref{tab:searchpath} records which taxonomy branch
was searched for which of our functions, and what it returned; the biologized
questions themselves are in Section~\ref{sec:biologized}.

\begin{table}[htbp]
\centering
\caption{Research path. Each of our functions was translated into a taxonomy
branch before any organism was considered.}
\label{tab:searchpath}
\renewcommand{\arraystretch}{1.25}
\small
\begin{tabularx}{\textwidth}{L{1.6cm} L{5.2cm} X}
\toprule
\textbf{Function} & \textbf{Taxonomy branch searched} & \textbf{Models it returned} \\
\midrule
F1, F5 & Modify $\rightarrow$ physical state $\rightarrow$ light/colour;
         Protect from non-living threats $\rightarrow$ light
       & Desert snail, \emph{Morpho}, \emph{Cyphochilus}, tree bark, cactus ribs \\
F2, F3 & Transform/convert energy $\rightarrow$ thermal, radiant;
         Distribute $\rightarrow$ energy
       & Elephant ears, elephant hairs, jackrabbit ears, Saharan silver ant \\
F4     & Store $\rightarrow$ energy
       & Sea star, desert snail \\
F6--F8 & Capture, absorb or filter $\rightarrow$ liquids; Store $\rightarrow$
         liquids; Distribute $\rightarrow$ liquids
       & Leaf stomata, camel fur, honeybee water collection, bryophytes \\
F9     & Sense signals $\rightarrow$ temperature \emph{paired with}
         Process signals $\rightarrow$ respond to signals
       & Honeybee fanning thresholds, leaf stomata, pine cone, elephant hot spots \\
F10, F11 & Protect from non-living threats $\rightarrow$ dirt/solids, light, ice
       & Lichen, lotus leaf \\
\bottomrule
\end{tabularx}
\end{table}

F9 is worth noting: it is the only one of our functions that does not sit in a
single branch of the taxonomy. It is the \emph{pairing} of \enquote{sense
signals: temperature} with \enquote{process signals: respond to signals}, which
is the same construction the Life's Principles checklist uses when it asks that
\enquote{the signal it sends, the antennae it uses to detect the signal, and its
response match} \parencite{baumeister2014}. That the central function of our
design question has to be assembled from two branches is itself a finding: it is
why searching by organism would not have produced it.

\paragraph{Sources, in the order they were used.} AskNature strategies and
innovations to find candidates \parencite{asknature}; the \emph{Genius of Biome}
report for the temperate broadleaf forest, which is the biome Zurich lies in,
for models already translated to the built environment \parencite{hok2013};
primary literature for the mechanisms; and commercial products where the
transfer has already been made, which are treated as state of the art rather
than as models. AskNature and AI-assisted search were used to \emph{find}
candidates only; every mechanism claim carried into this report is referenced to
a paper or to the organisation that published the strategy.

\paragraph{Which lenses were applied.} Of the five lenses in the method, two did
most of the work. The \emph{operating-condition lens} — intense solar radiation,
surfaces above \qty{60}{\celsius}, freeze--thaw, weeks without rain — led to the
desert organisms. The \emph{local lens} led to \parencite{hok2013}, because
Zurich sits in the temperate broadleaf forest biome and that report translates
exactly this biome into building design. The \emph{function lens} produced
Table~\ref{tab:searchpath}. The naturalist and ecological lenses contributed
context but no model that survived selection.

\subsubsection{Models found, and the selection decision}

\begin{table}[htbp]
\centering
\caption{Biological models found, and the selection or de-selection decision.
The criterion applied throughout: a model must be transferable to a
\emph{passive surface} that survives twenty years of the Swiss annual cycle.}
\label{tab:models}
\renewcommand{\arraystretch}{1.2}
\footnotesize
\begin{tabularx}{\textwidth}{L{2.9cm} L{1.3cm} X L{2.5cm}}
\toprule
\textbf{Organism / model} & \textbf{Function} & \textbf{Mechanism} & \textbf{Decision} \\
\midrule
\multicolumn{4}{l}{\textit{Prevention --- reduce what is absorbed}}\\
Desert snail \emph{Sphincterochila boissieri} & F1, F4 &
Shell reflects $\approx\qty{90}{\percent}$ of visible and
$\approx\qty{95}{\percent}$ of near-infrared light; in extreme heat the aperture
is sealed with a calcified plate and the animal withdraws, leaving an insulating
air space between body and hot ground \parencite{asknature_snail} &
\textbf{Selected.} Passive, mineral, non-metabolic --- the closest analogue to
an envelope we found \\
\emph{Morpho} butterfly & F1 &
Transparent chitin--air layers a few hundred nanometres apart diffract and
recombine light; colour comes from geometry, and the disordered structure keeps
the reflection diffuse rather than specular
\parencite{asknature_morpho} &
\textbf{Selected.} Decouples appearance from absorptance, which is what the
townscape constraint requires \\
White beetle \emph{Cyphochilus} & F1 &
Brilliant whiteness from a disordered chitin network only
$7\pm\qty{1.5}{\micro\metre}$ thick, without pigment \parencite{vukusic2007} &
\textbf{Selected} as the achromatic case of the same principle \\
Saharan silver ant \emph{Cataglyphis bombycina} & F1, F2 &
Triangular hairs scatter visible and near-infrared light and simultaneously
raise emissivity in the mid-infrared \parencite{shi2015} &
\textbf{Selected.} The only model that does F1 and F2 in one structure \\
Cactus ribs & F5, F3, F2 &
Peaks shade the troughs; the cooler trough air absorbs heat and rises out, while
the folds enlarge the emitting area \parencite{asknature_cactus} &
\textbf{Selected.} Geometry only --- no material, no water, nothing to bleach \\
Tree bark & F1, F2 &
Reflects part of the incoming radiation while emitting efficiently; rough
layering shades itself and traps insulating air \parencite{asknature_bark} &
\textbf{Selected.} A dead outer layer protecting a living structure --- the same
role our build-up has \\
\midrule
\multicolumn{4}{l}{\textit{Removal --- get heat out again}}\\
Elephant ears & F2, F3 &
Warm blood is carried from the core into large thin ears; the surface releases
it, and flapping raises the airflow \parencite{asknature_ears} &
\textbf{Selected.} Separates the place of absorption from the place of release \\
Elephant skin hairs & F3 &
Sparse hairs reach through the still boundary layer so moving air reaches the
skin; reported additional heat loss up to \qty{23}{\percent}
\parencite{asknature_leaves} &
\textbf{Selected.} Directly relevant at the low wind speeds of a canyon \\
Jackrabbit ears & F9, F2 &
Vasodilation varies the blood flow to a thin radiating surface; at
$\approx\qty{30}{\celsius}$ air temperature the ears shed all excess heat, and
the strategy exists to avoid spending water
\parencite{asknature_jackrabbit} &
\textbf{Selected.} The water-free route to a load-coupled response \\
Elephant skin hot spots & F9 &
Vasodilation in scattered patches rather than uniformly; the number of open
\enquote{thermal windows} varies with the load
\parencite{asknature_hotspots} &
\textbf{Selected.} Regulation by active \emph{area} rather than by material \\
\midrule
\multicolumn{4}{l}{\textit{Water path and storage}}\\
Leaf stomata & F9, F8 &
Guard cells swell and shrink with their water content and so open and close the
pore; evaporation itself pulls the water up, with no pump; leaf surfaces cool by
\qtyrange{11}{17}{\kelvin} \parencite[62--65]{hok2013} &
\textbf{Selected.} The clearest natural instance of F9 \\
Camel fur over sweat glands & F7, F8 &
Insulation over the evaporating layer; in the technical reproduction the pair
held a sample \qty{7}{\kelvin} below ambient for \qty{200}{\hour}, five times
longer than the wet layer alone \parencite{asknature_camel} &
\textbf{Selected.} Turns the water reserve from a quantity into a duration \\
Sea star & F4, F7, F9 &
Takes up more cold water at high tide after a hot low tide, so the previous
load sets the charge; fails once the sea is no longer colder than the air
\parencite{asknature_seastar} &
\textbf{Selected.} Charge-before-load, including its own failure mode \\
Bryophytes & F6, F7 &
No cuticle, no roots; a living layer over a decaying one holds moisture behind a
still-air boundary, and the plant survives drying out
\parencite[53--56]{hok2013} &
\textbf{Not carried} into the concepts --- superseded by the camel bilayer,
which achieves the same with a documented duration \\
\midrule
\multicolumn{4}{l}{\textit{Self-regulation and supporting functions}}\\
Honeybee fanning thresholds & F9 &
Each worker begins fanning at a slightly different temperature, and the spread is
genetic; the number of active bees rises with the load, so binary units produce
a graded colony response \parencite{asknature_bees} &
\textbf{Selected.} Resolves whether a threshold can satisfy F9 \\
Pine cone scales & F5, F9 &
Two bonded layers swell differently with moisture, so the scale bends; the cone
is dead tissue moved entirely by the weather
\parencite{asknature_pinecone} &
\textbf{Selected.} A passive actuator with no energy input \\
Lichen \emph{Lecanora conizaeoides} & F10, F11 &
Roughness at several scales keeps droplets perched so they roll off, while
hydrophobin-coated channels stay dry and admit air: waterproof and breathable at
once \parencite[41--44]{hok2013} &
\textbf{Selected.} A physical answer to soiling, not a chemical one \\
Desert shrub \emph{Encelia farinosa} & F1, F5 &
Dense leaf pubescence raises reflectance and is adjusted seasonally
\parencite{ehleringer1978} &
\textbf{Selected} as the biological precedent for a seasonally switching surface \\
African elephant skin cracks & F6--F8 &
A network of micrometre-scale cracks retains water and extends evaporative
cooling \parencite{martins2018} &
\textbf{Selected.} Third mechanism from the same organism \\
\midrule
\multicolumn{4}{l}{\textit{De-selected, with reasons}}\\
Termite mound & F3 &
The classic reference, but the \enquote{chimney air-conditioning} account is
contested; measurements support wind-driven transient gas exchange rather than a
thermosiphon \parencite{turner2008} &
\textbf{De-selected.} The evidence does not support the mechanism, and the
building usually cited for it uses fans \\
Namib desert beetle & F6 &
Hydrophilic--hydrophobic patterning harvests water from fog &
\textbf{De-selected.} Wrong season for our climate: fog in Zurich is a winter
phenomenon, and a heat wave brings neither rain nor fog \\
Thomson's gazelle carotid rete & F3 &
Counter-current heat exchange cools the brain &
\textbf{De-selected} as a whole-organism model --- internal physiology. The
counter-current \emph{arrangement} is retained as an option for the liquid loop
in Section~\ref{sec:creating} \\
Tenrec aestivation & --- &
Dormancy through the hot season &
\textbf{De-selected.} A behaviour; an envelope has none \\
Green façades and roofs & F1, F5 &
Shading and evapotranspiration from planting &
\textbf{Out of scope} by the project boundary (Table~\ref{tab:inout}), and in
direct competition with our own water reserve: planting needs most water exactly
when least is available \\
\bottomrule
\end{tabularx}
\end{table}

Two features of Table~\ref{tab:models} are deliberate. First, the same organism
appears three times --- the elephant supplies ear-based transport, boundary-layer
hairs and localised hot spots, and these are three different mechanisms serving
three different functions, not one model counted three times. Second, five
models were rejected \emph{with a stated criterion} rather than quietly dropped.
The criterion is the one at the head of the table, and it is worth stating
explicitly because it explains an asymmetry in what survived: most cooling in
nature is metabolic or behavioural, and an envelope has neither metabolism nor
behaviour. The two models that transfer most directly --- the snail shell and the
cactus rib --- are the two that are \emph{structures rather than processes}.

\todo{Before hand-in: open every AskNature page cited in this table and replace
it with the primary paper wherever one exists. Two are already identified: the
silver-ant film reproduction is \emph{Biologically inspired flexible photonic
films for efficient passive radiative cooling}, PNAS 2020, and the adaptive
infrared material is \emph{A dynamic thermoregulatory material inspired by squid
skin}, Nature Communications 2019. The course rule is explicit that AskNature
finds candidates and the paper is the evidence.}

\subsubsection{Mechanisms in depth}

Three mechanisms are developed further here, one per functional family. They
were chosen because each answers a question Section~2 left open.

\paragraph{Prevention --- why reflection alone is not the strategy.} The desert
snail survives on sun-exposed ground at up to \qty{50}{\celsius} without moving,
without water, and without a metabolism that could help. Its shell reflects
about \qty{90}{\percent} of visible and about \qty{95}{\percent} of near-infrared
light, which is above the $\alpha_\mathrm{sol}\geq0.80$ that KPI~2 asks for and
close to the ultra-white laboratory paints of \textcite{li2021} --- achieved with
calcium carbonate rather than a manufactured pigment. The part that matters for
us is what happens next. In extreme conditions the animal seals the aperture and
retracts into the second whorl, so an air space separates its tissue from the
shell wall nearest the hot ground. Reflection reduces what is absorbed; the air
space stops what is still absorbed from reaching what must be protected. For a
retrofit this maps onto a construction that already exists: a reflective outer
skin in front of a ventilated cavity. That matters for the night window, because
the \qty{7}{\kelvin} the city loses at night is heat released from mass that was
charged during the day, and a decoupled outer skin charges far less of it.

\paragraph{Regulation --- how a switch can behave like a controller.} In a
honeybee colony the temperature at which an individual worker begins to fan
varies between individuals, and the variation is genetic: a colony has one
queen but many fathers, so it contains patrilines with different thresholds. As
the nest warms, first a few bees fan, then more, then many, and the nest is held
below about \qty{36}{\celsius} with no central control at all. Each bee is a
binary switch; the colony response is continuous. This resolves an objection
that would otherwise disqualify every switching material considered in
Section~\ref{sec:creating}: a single threshold gives an almost infinite
$\partial P/\partial T_s$ at the switching point and roughly zero either side of
it, which is not the graded response KPI~1 asks for. A \emph{distribution} of
thresholds across many small units restores it. The same pattern appears in
\textcite{hok2013} as the \enquote{lots of littles} principle, arrived at
independently.

\paragraph{Water --- why metering matters more than harvesting.} A leaf cools
itself by \qtyrange{11}{17}{\kelvin} through pores whose two guard cells swell
and shrink with their water content, opening and closing the pore as they do;
the water is drawn up by the evaporation itself, with no pump anywhere
\parencite[62--65]{hok2013}. The valve responds directly to heat and dryness, so
the effort tracks the load. The camel supplies the complementary half. Its fur
is insulation worn \emph{in} the heat, and in the technical reproduction of that
arrangement --- an evaporating hydrogel underneath a heat-blocking but
vapour-permeable aerogel --- a sample was held \qty{7}{\kelvin} below ambient for
\qty{200}{\hour}, five times longer than the wet layer alone
\parencite{asknature_camel}. Insulating the wet layer, rather than exposing it,
is what converts a fixed volume of water into a duration. Since KPI~6 asks for
five days of autonomy and Zurich supplies no rain during the design heat wave,
metering is the binding problem and harvesting is not.

\todo{Add one figure per mechanism above. Prompts for generating them are in
\texttt{Working documents/image\_prompts.md}; a generated figure is an
illustration and must be captioned as schematic, and the AI use belongs in
Appendix~A.}

\subsubsection{Abstraction into design principles}

Table~\ref{tab:abstraction} keeps the three columns strictly separate, as the
method requires: what was observed, what the observation abstracts to, and what
it might become technically. Only the third column contains anything that could
be built, and nothing in it has been selected yet --- that is the work of
Section~\ref{sec:creating}.

\begin{table}[htbp]
\centering
\caption{From biological observation to transferable design principle. The third
column lists possibilities, not decisions.}
\label{tab:abstraction}
\renewcommand{\arraystretch}{1.25}
\footnotesize
\begin{tabularx}{\textwidth}{X X X}
\toprule
\textbf{Biological observation} & \textbf{Abstracted principle} & \textbf{Possible technical interpretation} \\
\midrule
The snail's shell reflects strongly \emph{and} an air space separates it from the
hot substrate
& Reflect first, then decouple whatever still gets through with a layer of air
& Reflective outer skin over a ventilated cavity; structural white instead of a
pigment \\
\emph{Morpho} colour comes from layer geometry, not pigment, and the disordered
structure stays diffuse
& Separate what a surface looks like from what it does with heat
& Structurally coloured or near-infrared-selective finish that keeps the
existing façade colour \\
Cactus peaks shade the troughs; cooler trough air absorbs heat and rises out
& Use folded geometry for self-shading, more emitting area and better airflow
& Ribbed or profiled façade and paving relief \\
Sparse hairs reach through the still air layer clinging to the skin
& Break the boundary layer so moving air reaches the surface
& Millimetre-scale texture on façade and pavement \\
The jackrabbit varies blood flow to a thin radiator; the elephant opens patches
& Regulate by varying the \emph{transport} to a radiating surface, and place the
radiator only where conditions are best
& Pumpless density-driven liquid loop from the canyon face to a parapet-level
radiator with a good sky view \\
Guard cells open and close the pore in response to heat and dryness
& Build the control into the opening, so the escape route widens with the load
& Load-driven valve or wick, rather than a controlled one \\
Camel fur insulates \emph{over} the evaporating skin
& Cut the load first, then spend a little water on the remainder, and cover the
wet layer so the reserve lasts
& Absorbent layer beneath a vapour-permeable cover \\
Worker bees each switch at a different temperature
& Spread many binary elements across a range of trigger points and the sum
responds continuously
& Any switching element built as many small units with a deliberate spread of
switching temperatures \\
Pine cone scales bend because two bonded layers swell differently
& Make moving parts from two layers that respond differently and let the
environment actuate them
& Bilayer shading element; hygroscopic or thermal actuation, no motor \\
The sea star charges on cold water sized by the previous hot exposure
& Charge the store during the cool phase, sized by the load of the previous cycle
& Night-charged store, with an explicit reset condition across the five-day case \\
Lichen roughness at several scales sheds water while staying breathable
& Use geometry rather than chemistry to shed water and stay clean
& Multi-scale textured finish; no biocide in the film \\
\bottomrule
\end{tabularx}
\end{table}

\paragraph{One principle, found three times.} The jackrabbit and the elephant
regulate by how much \emph{area} is active, the bees by how many \emph{units}
have switched, and a bilayer shading element by how far it has \emph{deployed}.
All three modulate how many elements act rather than how strongly each one acts,
and all three arrive at a graded response without a controller. That convergence
from three unrelated organisms is the strongest abstraction this project has
produced, and it is the one carried into every concept in
Section~\ref{sec:creating}.

\subsubsection{Life's Principles evident in the selected strategies}

The principles below are the ones the \emph{selected models} actually exhibit;
Section~\ref{sec:sustainability} assesses how well our final concept applies
them, which is a different question.

\begin{itemize}
  \item \textbf{Be locally attuned and responsive} --- present in every model
        that couples its response to the local load, and named by the checklist
        in the form our F9 requires: the signal, the antenna that detects it and
        the response must match \parencite{baumeister2014}.
  \item \textbf{Be resource efficient (material and energy)}, sub-principle
        \emph{use low energy processes} --- none of the selected mechanisms
        consumes energy; the load itself drives them.
  \item \textbf{Adapt to changing conditions}, sub-principles \emph{incorporate
        diversity} and \emph{maintain integrity through self-renewal} ---
        \emph{Encelia} and the pine cone for the first, the lichen for the
        second.
  \item \textbf{Use life-friendly chemistry} --- the lichen and the
        \emph{Cyphochilus} structure both achieve with geometry what industry
        achieves with a compound that then has to be justified.
\end{itemize}

\todo{Section~4 must \emph{assess} these, not list them again. The rubric marks
down a list.}

%-----------------------------------------------------------------------
\subsection{Creating: Engineering Concept Development}
\label{sec:creating}

\subsubsection{Ideation}

\todo{Run and document Crazy~8s: eight technically different ways of
transferring the abstracted principles of Table~\ref{tab:abstraction}. Vary the
mechanism, not the shape, and do not evaluate while generating. Include a photo
or scan of the sheets. Warning: four of the eight ideas in the lecture's own
worked example (night flush, evaporative inlet, thermal buffer, shading skin)
overlap our function list. If our eight reproduce that list we have copied the
example instead of using our own abstractions.}

\subsubsection{Morphological matrix}

The rows are chosen, not inherited. Of our thirteen functions, a row belongs in
the matrix only where alternative principles would lead to a meaningfully
different concept. On that criterion F0 is excluded as the main function, F10 and
F11 are requirements and engineering checks rather than sources of concept
variants, and F12 is included only if the mounting really changes the
architecture. That leaves five rows.

\begin{table}[htbp]
\centering
\caption{Morphological matrix. Rows are the five functions for which
alternative principles produce genuinely different concepts; columns are the
technical interpretations abstracted in Table~\ref{tab:abstraction}.}
\label{tab:morph}
\renewcommand{\arraystretch}{1.3}
\small
\begin{tabularx}{\textwidth}{L{3.4cm} X X X}
\toprule
\textbf{Function} & \textbf{Option A} & \textbf{Option B} & \textbf{Option C} \\
\midrule
F1, F5 \newline Reduce heat input
& Structural white or near-infrared-selective finish
& Ribbed, self-shading geometry
& Deployable shading element \\
F2, F3 \newline Move heat away
& Broadband emitter with a good sky view
& Boundary-layer texture for convection
& Density-driven liquid loop to a parapet radiator \\
F4 \newline Store and buffer
& Thermal mass in the build-up
& Latent store (phase change)
& No added storage; decoupling only \\
F6--F8 \newline Handle water
& Absorbent layer under a vapour-permeable cover
& Wick with load-driven release
& No water path \\
F9 \newline Couple response to load
& Switching material, many units, spread thresholds
& Load-driven valve or pore
& Bilayer element that deploys with heat or moisture \\
\bottomrule
\end{tabularx}
\end{table}

\todo{Read at least three coherent combinations off this matrix and name them.
A column is not a concept: check that the options combined are physically
compatible. Then select two that are technically distinct --- different physics,
not different styling --- and develop them below.}

\subsubsection{Concept A --- \todo{name}}
\todo{Annotated drawing plus description: components, functional chain, flows,
candidate materials, key technical parameters. Annotate inlet, channels, driving
effect, outlet, heat path, materials, key dimensions and expected performance,
and state explicitly which parts are inspired by biology and which are
conventional engineering.}

\subsubsection{Concept B --- \todo{name}}
\todo{Must be technically distinct from A. The four directions currently on the
table are: a switching coating (prevention, no night effect); a self-deploying
shade (prevention plus an open night sky, acts at pedestrian level); an
evaporative build-up with a water path (removal, metered per the camel bilayer);
and a ribbed geometry (prevention and convection, no material and no water). The
pairing with the widest physical separation is the deploying shade against the
evaporative build-up.}

\subsubsection{Engineering Canvas and critical questions}

\begin{table}[htbp]
\centering
\caption{Engineering Canvas. One per concept; the output of the canvas is the
list of critical questions below it.}
\label{tab:canvas}
\renewcommand{\arraystretch}{1.3}
\small
\begin{tabularx}{\textwidth}{L{4.2cm} X}
\toprule
\textbf{Field} & \textbf{Concept \todo{A / B}} \\
\midrule
Function --- what must it achieve? & \todo{} \\
Biological mechanism & \todo{} \\
Abstracted principle & \todo{} \\
Technical principle --- what is transferred? & \todo{} \\
Architecture --- components and interfaces & \todo{} \\
Inputs / outputs --- material, energy, information & \todo{} \\
Requirements / KPIs & \todo{} \\
Assumptions & \todo{} \\
Uncertainties --- what could invalidate it? & \todo{} \\
Next check --- what should we calculate or test? & \todo{} \\
\bottomrule
\end{tabularx}
\end{table}

A critical question names a quantity, a threshold and a scale. Three follow
directly from the abstractions of Section~\ref{sec:discovering} and are carried
forward as candidates:

\begin{enumerate}
  \item Is the buoyancy in a pumpless density-driven loop sufficient to move the
        required heat flux over the height available in a five-storey block, and
        does the loop survive a Swiss winter?
  \item At what wind speed must a self-deploying shading element furl, and what
        furls it? A thermal actuator cannot sense wind, so an independent and
        equally passive release is required before the concept is safe over a
        street.
  \item Can a switching element be produced with a deliberate \emph{spread} of
        switching temperatures, as the honeybee threshold model requires, or does
        the chemistry force a single sharp threshold?
  \item \todo{Select two or three of these, or replace them once the concepts are
        fixed. They must be the assumptions that could change the design
        decision.}
\end{enumerate}

\subsubsection{Engineering check 1: steady-state surface energy balance}

\todo{This check is prepared because it is the most direct way to test KPI~1 and
KPI~3. Fill in your own numbers, state every assumption, and keep the units.}

For an opaque surface in steady state, the energy balance per unit area is

\begin{equation}
(1-\alpha_\mathrm{sol})\,G \;+\; \varepsilon L_\downarrow
\;-\; \varepsilon \sigma T_s^{4} \;-\; h_c\,(T_s - T_a) \;-\; q_\mathrm{cond}
\;=\; 0 ,
\label{eq:energybalance}
\end{equation}

\noindent where $\alpha_\mathrm{sol}$ is the solar reflectance,
$G$ the global irradiance on the surface, $\varepsilon$ the thermal emittance,
$L_\downarrow$ the incoming long-wave radiation from the sky,
$\sigma = \qty{5.67e-8}{\watt\per\square\metre\per\kelvin\tothe{4}}$ the
Stefan--Boltzmann constant, $T_s$ the surface temperature, $T_a$ the air
temperature, $h_c$ the convective heat transfer coefficient and
$q_\mathrm{cond}$ the heat conducted into the construction.

\begin{table}[htbp]
\centering
\caption{Assumptions for the energy balance. Every value needs a source.}
\label{tab:assumptions}
\renewcommand{\arraystretch}{1.2}
\small
\begin{tabularx}{\textwidth}{L{3.6cm} L{3.0cm} X}
\toprule
\textbf{Quantity} & \textbf{Value used} & \textbf{Justification / source} \\
\midrule
$G$              & \todo{} \unit{\watt\per\square\metre} & \todo{MeteoSwiss, clear-sky summer day; Table~\ref{tab:windows} gives \qtyrange{900}{950}{\watt\per\square\metre} horizontal} \\
$T_a$            & \todo{} \unit{\celsius} & \todo{Table~\ref{tab:windows}: \qtyrange{33}{36}{\celsius} in the load window} \\
$h_c$            & \todo{} \unit{\watt\per\square\metre\per\kelvin} & \todo{low wind speed in a street canyon} \\
$L_\downarrow$   & \todo{} \unit{\watt\per\square\metre} & \todo{clear-sky estimate; state the correlation used} \\
$\alpha_\mathrm{sol}$ reference & $0.3$ & Value the city uses for un-greened façades \parencite[114]{stadtzuerich2020} \\
$\alpha_\mathrm{sol}$ concept   & \todo{} & \todo{target $\geq0.80$, KPI~2} \\
$\varepsilon$    & \todo{} & \todo{state whether the concept changes emittance or only reflectance --- these are separate questions} \\
\bottomrule
\end{tabularx}
\end{table}

\paragraph{Result and meaning.} \todo{Solve Eq.~\eqref{eq:energybalance} for
$T_s$ for the reference and for the concept, report both, and state explicitly
what the difference does and does \emph{not} prove (see the note on KPI~3 in
Section~\ref{sec:kpi}).}

\subsubsection{Engineering check 2: \todo{title}}

\todo{Options that are realistic within the week: (a) a simple outdoor
measurement --- identical plates with different surfaces, logged over a few
hours; (b) a transient 1-D model of the build-up over a diurnal cycle;
(c) a water balance --- litres of rainwater per square metre per hot day against
the latent heat of vaporisation of \qty{2.45}{\mega\joule\per\kilo\gram},
checked against the \qty{200}{\hour} the camel bilayer achieved
\parencite{asknature_camel}; (d) a material and cost comparison. Report method,
raw data, result and uncertainty.}

\paragraph{A scale check that applies to four of our models.} Micro-scale
effects do not automatically survive being made square metres large. The
\emph{Cyphochilus} whiteness, the \emph{Morpho} colour, the lichen roughness and
the stomatal valve all live at micrometre scale, and the failure mode is named
explicitly in the method: upscaling can break the function. \todo{State for each
of the four whether the effect scales, and cite the evidence. For the lotus
effect the answer is partly known, because a lotus-effect façade paint has been
manufactured since 1999 \parencite{sto_lotusan}.}

%-----------------------------------------------------------------------
\subsection{Concept Evaluation \& Iteration}
\label{sec:evaluation}

\begin{table}[htbp]
\centering
\caption{Decision matrix. Ratings 1--5, weights from
Section~\ref{sec:kpi}. Every rating must be justified in the text by a
calculation or a reasoned estimate.}
\label{tab:decision}
\renewcommand{\arraystretch}{1.25}
\small
\begin{tabularx}{\textwidth}{L{5.2cm} c c c X}
\toprule
\textbf{Criterion} & \textbf{Weight} & \textbf{Concept A} & \textbf{Concept B} & \textbf{Evidence used} \\
\midrule
KPI~1 -- Self-regulation gradient   & \todo{} & \todo{} & \todo{} & \todo{} \\
KPI~2 -- Solar reflectance          & \todo{} & \todo{} & \todo{} & \todo{} \\
KPI~3 -- Peak surface temperature   & \todo{} & \todo{} & \todo{} & \todo{} \\
KPI~4 -- PET at \qty{2}{\metre}     & \todo{} & \todo{} & \todo{} & \todo{} \\
KPI~5 -- Night radiative cooling    & \todo{} & \todo{} & \todo{} & \todo{} \\
KPI~6 -- Water path and autonomy    & \todo{} & \todo{} & \todo{} & \todo{} \\
KPI~7 -- Durability of the function & \todo{} & \todo{} & \todo{} & \todo{} \\
\midrule
Feasibility                          & \todo{} & \todo{} & \todo{} & \todo{} \\
Quality of the biomimetic transfer   & \todo{} & \todo{} & \todo{} & \todo{} \\
Evidence and uncertainty             & \todo{} & \todo{} & \todo{} & \todo{} \\
\midrule
\textbf{Weighted total}              &         & \todo{} & \todo{} & \\
\bottomrule
\end{tabularx}
\end{table}

The last three rows are not KPIs and are deliberately separate from them. They
carry the questions a KPI table cannot answer: whether the concept can be built,
whether it genuinely transfers the biology or merely resembles it, and how well
supported the ratings above are. The third of these is the one that prevents a
confident-looking table from concealing guesswork.

\subsubsection{Review and feedback}
\todo{World café peer feedback, instructor feedback, and AI used as a critical
reviewer. For every AI comment record the decision --- accept, reject or
investigate --- and the justification. A documented rejection is worth more than
a list of accepted comments, because it shows judgement. This belongs in
Appendix~A as well.}

\subsubsection{Selection and Final Concept V2}

\todo{State which concept, or which justified hybrid, was selected, and describe
V2 in detail. The minimum evidence of iteration is one named weakness or
uncertainty in the earlier concept and a documented response to it in V2.}

\begin{table}[htbp]
\centering
\caption{Final output of the creating and evaluating phase:
Concept V1a/b $\rightarrow$ engineering review $\rightarrow$ evidence
$\rightarrow$ Concept V2.}
\label{tab:v1v2}
\renewcommand{\arraystretch}{1.25}
\small
\begin{tabularx}{\textwidth}{L{4.0cm} X X}
\toprule
\textbf{What changed} & \textbf{Why it changed} & \textbf{Evidence that informed the change} \\
\midrule
\todo{} & \todo{} & \todo{} \\
\todo{} & \todo{} & \todo{} \\
\bottomrule
\end{tabularx}
\end{table}

\subsubsection{Remaining uncertainties}
\todo{What is still unproven, and what should be tested or developed next. The
scale question above and the wind release for a deploying element are the two
that are already known to be open.}
